%% This document created by Scientific Word (R)

\documentstyle[12pt,qqaajart]{article}

\input tcilatex

\QQQ{Language}{
American English
}

\begin{document}


\begin{center}
{\Large A Remark on Bartels and Conn's Linearly Constrained Discrete $l_1$ Problems}\footnote{Financial support from the National
Science Foundation Grants SES 89-22472 is gratefully acknowledged.
Authors' addresses: Roger W. Koenker, Department of Economics, University
of Illinois, Champaign, IL 61820.  Pin T. Ng, Department of Economics,
University of Houston, Houston, TX 77204-5882}

\bigskip\ 

{\normalsize Roger W. Koenker}

{\normalsize Department of Economics }

{\normalsize University of Illinois at Urbana-Champaign }

{\normalsize and }

{\normalsize Pin T. Ng }

{\normalsize Department of Economics }

{\normalsize University of Houston }

\bigskip\ 
\end{center}
\noindent
{\it Categories and Subject Descriptors:} G.1.6 [Numerical Analysis]: 
Optimization -- constrained optimization, gradient methods, linear 
programming; G.3 [Mathematics of Computing]: Probability and Statistics --
statistical computing\\
{\it General Terms:} Algorithms, Performance, Reliability, Theory\\
{\it Key Words and Phrases:} degeneracy, $l_1$ norm, discrete approximation,
linearly constrained approximation \\

\begin{center}
{\bf Abstract}
\end{center}

Two modifications of Bartels and Conn's algorithm for solving linearly
constrained discrete $l_1$ problems are described.  The modifications
are designed to improve performance of the algorithm under conditions 
of degeneracy.

\section{Introduction}

Bartels and Conn \cite{BC:1980} provided an efficient algorithm to solve the
following linearly constrained discrete $l_1$ problem, 
\begin{equation}
\label{eqn:bcobj1}\min _{x\in R^n}\ \ ||A^{\prime }x-b||_1 
\end{equation}
$$
\mbox{subject to}\ \ \ \ G^{\prime }x\ =\ h,\ \ \ \ C^{\prime }x\ \geq \ d. 
$$
where $A,G,C$ are matrices with columns $a_i,g_j,c_k\in R^n$ respectively,
and $b,h,d$ are vectors with elements $\beta _i,\eta _j,\delta _k$
respectively. Letting $E=\left[ e_{1,}\cdots ,e_m\right] =\left[ \gamma
A\vdots G\vdots C\right]$ and $f=\left[ \zeta
_1,\cdots ,\zeta _m\right] ^{\prime }=\left[ \gamma b^{\prime }\vdots
h^{\prime }\vdots d^{\prime }\right] ^{\prime }\in {\bf R}^m$, (\ref
{eqn:bcobj1}) can be written as 
\begin{equation}
\label{eqn:bcobj2}\min _{x\in R^n}\psi \left( x\right) =\sum_\mu ^{}|e_\mu
^{\prime }x\ -\ \zeta _\mu |\ -\ \sum_\nu ^{}\min (0,e_\nu ^{\prime }x\ -\
\zeta _\nu ) 
\end{equation}
where $\gamma $ is a Lagrange multiplier, $\mu $ is the index for the
columns of $\left[ \gamma A\vdots G\right] \,$ and $\nu $ is the index for
the columns of $\left[ C\right] $.

Our motivation for studying the algorithm is that it is well suited for
computing the quantile smoothing splines described in Koenker, Ng, and
Portnoy \cite{KNP:1994}.  However, in the course of adapting the algorithm for
this purpose we encountered some difficulties with the behavior of the
algorithm in degenerate situations.  Since our problems are extremely
sparse and degeneracy is the norm rather than an exceptional situation, 
it proved essential to understand and correct this anomalous behavior.
We describe these modifications briefly in the next section; a detailed
description of the quantile smoothing spline computations will be
reported elsewhere.  Throughout, we adhere as closely as possible to the
notation introduced by Bartels and Conn.


\section{Degeneracy in the Bartels and Conn  Algorithm}

The directional derivative of (\ref{eqn:bcobj2})
is written by Bartels and Conn as
\begin{equation}
\label{eqn:dd}
\psi^{\prime }\left( x, p\right)  = 
g^{\prime }p+ \left[ \sum_{\mu \in {\cal E}}\left| e_\mu
^{\prime }p\right| -\sum_{\nu \in {\cal E}}\min \left( 0,e_\nu ^{\prime
}p\right) \right] 
\end{equation}
where the {\it restricted gradient} is 
\begin{equation}
\label{eqn:gradient}g\ =\ \sum_{\mu \NEG \in {\cal E}}^{}\sigma _\mu e_\mu
+\sum_{\nu \NEG \in {\cal E}}^{}\sigma _\nu ^{-}e_\nu 
\end{equation}
\begin{equation}
\label{eqn:sigma} 
\begin{array}{lll}
\sigma _\mu & = & \left\{ 
\begin{array}{lllll}
1 & if & e_\mu ^{\prime }x-\zeta _\mu & > & 0 \\ 
-1 & if & e_\mu ^{^{\prime }}x-\zeta _\mu & < & 0 
\end{array}
\right. \\ 
\sigma _\nu ^{-} & = & \left\{ 
\begin{array}{lllll}
0 & if & e_\nu ^{\prime }x-\zeta _\nu & > & 0 \\ 
-1 & if & e_\nu ^{^{\prime }}x-\zeta _\nu & < & 0 
\end{array}
\right. 
\end{array}
\end{equation}
and $\cal E$ denotes the index set of the zero elements of the residual
vector $E^{\prime}x-f$.  A minimal set of linearly independent columns 
$\{ e_i: \; i \in \cal E \}$ are referred to as the active columns.
Degeneracy occurs when the full set of zero residual columns are
linearly dependent.  When this occurs, Bartels and Conn's strategy 
is to randomize the corresponding $\sigma$'s 
in  ({\ref{eqn:gradient}) inside the \{ {\bf find p }\} module. 
A ``descent direction'' is then constructed from this randomized restricted 
gradient.

The problem, as we see it, with this strategy is that the resulting 
restricted gradient
will not accurately reflect the contributions of the non-active, 
zero-residual columns to the directional derivitive.  
{\it These contributions depend on the chosen direction $p$}.  
Ignoring this can result in the algorithm taking steps in 
directions which actually increase, rather
than decrease, the objective function.  To correct this, we suggest
computing the direction $p$ as in Bartels and Conn, 
decomposing the index set of  non-active columns into a zero residual set
$\overline { \cal {E}}_0$ and  the others
$\overline { \cal {E}}_1$, recomputing the restricted gradient as
\begin{equation}
\label{eqn:gradient1}
g(p) \ =\ \sum_{\mu \in \overline{{\cal E}}_0}^{}\widehat{%
\sigma }_\mu e_\mu +\sum_{\mu \in \overline{{\cal E}}_1}^{}\sigma _\mu e_\mu
+\sum_{\nu \in \overline{{\cal E}}_0}^{}\widehat{\sigma }_\nu ^{-}e_\nu
+\sum_{\nu \in \overline{{\cal E}}_1}^{}\sigma _\nu ^{-}e_\nu 
\end{equation}
where 
\begin{equation}
\label{eqn:sigmahat} 
\begin{array}{lll}
\widehat{\sigma }_\mu & = & \left\{ 
\begin{array}{lllll}
1 & if & e_\mu ^{\prime }p & > & 0 \\ 
0 & if & e_\mu ^{\prime }p & = & 0 \\ 
-1 & if & e_\mu ^{^{\prime }}p & < & 0 
\end{array}
\right. \\ 
\widehat{\sigma }_\nu ^{-} & = & \left\{ 
\begin{array}{lllll}
0 & if & e_\nu ^{\prime }p & \geq & 0 \\ 
-1 & if & e_\nu ^{^{\prime }}p & < & 0 .
\end{array}
\right. 
\end{array}
\end{equation}
and finally, checking whether  $p^{\prime}g(p) < 0$ before proceeding 
(in \{ {\bf {find $\alpha$}}\}) to take a step in direction $p$.
A further source of difficulties is also apparent from ({\ref{eqn:sigmahat}}).
Since $e_\mu^{\prime} p$ may equal zero as well as $\pm 1$
it is important in the randomization step to randomize among all {\it three}
possible values.  If this is not done, the algorithm can cycle
indefinitely never finding a direction of descent.  With somewhat more
effort, one could further modify the algorithm to implement a
deterministic strategy for handling the problem of degeneracy as
discussed for example in Gill, Murray and Wright (1991, Section 8.3.3).

Having made these two rather minor modifications:  randomizing over
$\{ -1,0,1 \}$ rather than $\{ -1,1\}$, and verifying under conditions
of degeneracy that the randomized direction $p$ is indeed a descent
direction by computing an updated restricted gradient -- we find that
the Bartels and Conn algorithm performs extremely well for our quantile
spline applications.



\begin{thebibliography}{99}
\bibitem{BC:1980}  Bartels, R. and Conn, A. 
Linearly constrained discrete $%
l_1$ problems, {\em ACM Trans. Math. Softw., 6}, 4, (1980), 
594-608.

\bibitem{KNP:1994}  Koenker, R., Ng, P. and Portnoy, S. L. 
Quantile smoothing splines, {\em Biometrika, 84}, 4, (1994), 673-680.

\bibitem{GMW:1991}  Gill, P.E., Murray, W., and Wright, M.H.
{\it Numerical Linear Algebra and Optimization}, 1991, Addison-Wesley.
\end{thebibliography}

\end{document}
